\section{高阶导数}
\label{sec:ru-higher-derivatives}

导函数本身也可以求导。若$f'$在某点可导，它在该点的导数称为$f$的二阶导数，记作$f''$；依次定义$f^{(n)}=(f^{(n-1)})'$，约定$f^{(0)}=f$。
例如$f(x)=x^3-3x$时，$f'(x)=3x^2-3$，$f''(x)=6x$，$f'''\equiv6$，四阶及以上导数为零；$e^x$的各阶导数都是$e^x$；$\sin x$的导数按$\cos x,-\sin x,-\cos x,\sin x$四步一循环。

高阶导数存在与高阶导数连续是两个不同的条件。若$f$在开区间$I$上直到$n$阶导数都存在且连续，称$f$为$n$阶\term{连续可微}（continuously differentiable）函数，记作$f\in C^n(I)$；各阶导数都存在时记作$f\in C^\infty(I)$。
多项式、指数和三角函数都属于$C^\infty$。下面的例子说明处处可导的函数，导函数可以不连续。
\begin{example}[处处可导而导函数不连续]
\label{ex:analysis-derivative-not-continuous}
令$f(0)=0$，$x\neq0$时$f(x)=x^2\sin(1/x)$。在原点，差商$h\sin(1/h)$被$\pm|h|$夹住，趋于$0$，所以$f'(0)=0$。在$x\neq0$处，由乘积法则和链式法则，
\[
 f'(x)=2x\sin\frac1x-\cos\frac1x.
\]
取$x_n=1/(2n\pi)\to0$，则$f'(x_n)=-1$，不趋于$f'(0)=0$。所以$f$处处可导，但不属于$C^1(\mathbb R)$，更没有二阶导数$f''(0)$。
\end{example}
导数在原点附近的振荡来自$-\cos(1/x)$一项。可见“$n$阶可导”与“属于$C^n$”需要分别验证；后面的Taylor公式会明确写出用到的是哪一种条件。

乘积的高阶导数有类似二项式定理的公式。对$n$归纳，并用组合数的递推$\binom nk+\binom n{k-1}=\binom{n+1}k$合并同类项，得到Leibniz公式
\[
 (fg)^{(n)}=\sum_{k=0}^n\binom nkf^{(k)}g^{(n-k)},
\]
其中$f,g$在所讨论的点$n$阶可导。例如$(x^2e^x)^{(n)}$中，$x^2$三阶以上的导数为零，只剩三项：
\[
 (x^2e^x)^{(n)}=x^2e^x+n\cdot2x\,e^x+\binom n2\cdot2\,e^x=e^x\bigl(x^2+2nx+n(n-1)\bigr).
\]

复合函数逐层求导时，每次求导都要再用一次链式法则或乘积法则。以第~\ref{sec:ru-function-arithmetic}节的复合函数$g(x)=\sqrt{1-x^2}$为例，在$(-1,1)$上，
\[
 g'(x)=\frac{-2x}{2\sqrt{1-x^2}}=-\frac x{\sqrt{1-x^2}},\qquad
 g''(x)=-\frac{\sqrt{1-x^2}+x^2/\sqrt{1-x^2}}{1-x^2}=-\frac1{(1-x^2)^{3/2}}.
\]
对参数函数族，求导的变量要说清楚：$f_{a,b}(x)=ax+b$关于输入$x$的导数是$a$，二阶导数为$0$；若把$x$固定、改变参数$a$，函数值的变化率则是$x$。
这些逐层计算得到的各阶导数，正是下面Taylor展开所需的系数。
